\documentclass[11pt,a4paper]{article}
\usepackage[utf8]{inputenc}
\usepackage[T1]{fontenc}
\usepackage[brazilian]{babel}
\usepackage[margin=2cm]{geometry}
\usepackage{lmodern}
\usepackage{amssymb}
\usepackage{xcolor}
\usepackage{array}
\usepackage{booktabs}
\usepackage{enumitem}
\usepackage{titlesec}
\usepackage{fancyhdr}
\usepackage{parskip}

\definecolor{brandgreen}{HTML}{2E5E3E}
\definecolor{lightbg}{HTML}{F5F2EC}

\titleformat{\section}{\large\bfseries\color{brandgreen}}{}{0pt}{}[{\color{brandgreen}\titlerule}]
\titleformat{\subsection}{\bfseries\color{brandgreen}}{}{0pt}{}

\pagestyle{fancy}
\fancyhf{}
\fancyhead[L]{\small Querência ERP --- Avaliação de usabilidade}
\fancyhead[R]{\small Roteiro --- Papel Associado}
\fancyfoot[C]{\small\thepage}
\renewcommand{\headrulewidth}{0.4pt}

\newcommand{\dizer}[1]{%
  \noindent\colorbox{lightbg}{%
    \begin{minipage}{0.95\linewidth}
    \vspace{1mm}
    \textbf{O que dizer ao participante:} ``#1''
    \vspace{1mm}
    \end{minipage}%
  }\vspace{3mm}
}

\newcommand{\criterio}[1]{%
  \noindent\textit{Critério de sucesso objetivo:} #1
}

\newcommand{\estado}[1]{%
  \noindent\textit{Estado necessário (já preparado no kit):} #1
  \vspace{2mm}
}

\begin{document}

\begin{center}
  {\LARGE\bfseries\color{brandgreen} Roteiro do Participante}\\[2mm]
  {\large Papel: Associado (app mobile)}
\end{center}

\vspace{2mm}
\noindent\colorbox{lightbg}{%
  \begin{minipage}{0.95\linewidth}
  \vspace{1mm}
  Este participante executa os três fluxos abaixo \textbf{em sequência}, sempre no aplicativo
  mobile, no celular do moderador. Nenhuma tarefa depende da outra nem de outro participante ---
  todo o estado necessário já foi preparado no kit (ver \texttt{roteiro-moderador.pdf}). Entre um
  fluxo e outro, o \textbf{moderador} troca de conta no app (não é tarefa avaliada).
  \vspace{1mm}
  \end{minipage}%
}

\section*{Fluxo 1 --- Criar acesso como associado (RF01, RF18)}

Fazer \textbf{uma} das duas variantes abaixo (combinada previamente com o moderador).

\subsection*{Variante A --- Auto-cadastro}

\estado{ficha de persona ainda não usada.}

\dizer{Você quer se tornar associado da entidade e está criando sua conta pelo aplicativo pela
primeira vez.}

\begin{enumerate}[leftmargin=6mm]
  \item Partir da tela de login (app deslogado).
  \item Tocar \textbf{``Ainda não é associado? Cadastre-se''}.
  \item Preencher com os dados da ficha de persona do kit: Nome, CPF, Telefone (seção ``Dados
    pessoais''), Endereço (seção com ícone de mapa), E-mail e Senha.
  \item Tocar \textbf{Criar conta}.
  \item Esperado: mensagem de sucesso / tela indicando ``pendente de validação''.
\end{enumerate}

\textbf{Verificação objetiva:}
\begin{verbatim}
cd backend
node scripts/verificar-tarefas.js 1 <CPF-DA-FICHA>
\end{verbatim}
\criterio{associado existe com status \texttt{pendente\_validacao} \textbf{e} conta de login
vinculada.}

\subsection*{Variante B --- Vincular conta}

\estado{CPF pré-cadastrado pelo moderador (ativo, sem login), anotado no cartão do kit.}

\dizer{A diretoria já cadastrou você como associado, mas você nunca entrou no aplicativo. Use o
CPF que consta neste cartão para criar seu acesso.}

\begin{enumerate}[leftmargin=6mm]
  \item Partir da tela de login (app deslogado).
  \item Tocar \textbf{``Já é associado mas nunca acessou?''}.
  \item Preencher o CPF do cartão, e-mail e senha novos.
  \item Tocar \textbf{Vincular conta}.
\end{enumerate}

\textbf{Verificação objetiva:}
\begin{verbatim}
node scripts/verificar-tarefas.js 1 <CPF-DO-CARTAO>
\end{verbatim}
\criterio{ambos os critérios (status \texttt{ativo} + conta de login vinculada).}

\vspace{1mm}
\noindent\textit{Sem encadeamento: a conta criada aqui não é usada nos Fluxos 2 e 3. O moderador
sai dela e entra na conta de teste do kit antes de seguir.}

\section*{Fluxo 2 --- Compra de ingresso (RF12, RF20, RF21)}

\estado{conta de teste do kit, já logada pelo moderador; evento publicado com ingresso avulso
configurado.}

\begin{enumerate}[leftmargin=6mm]
  \item Abrir a aba \textbf{Eventos}, tocar no evento indicado no kit.
  \item Na tela do evento, tocar \textbf{Comprar ingressos}.
  \item Se aparecer o seletor ``Mesa / Ingresso avulso'', escolher \textbf{Ingresso avulso} (a
    menos que o moderador peça explicitamente a variante de reserva de mesa).
  \item Confirmar a compra do próprio ingresso (e, opcionalmente, adicionar 1 acompanhante).
  \item Abrir a aba \textbf{Ingressos}, localizar o ingresso recém-comprado --- confirmar que
    aparece um código (\texttt{\#001}, \texttt{\#002}\ldots) junto ao nome, e o QR ao tocar no
    card.
\end{enumerate}

\textbf{Verificação objetiva:}
\begin{small}
\begin{verbatim}
node scripts/verificar-tarefas.js 2 <CPF-DA-CONTA-DE-TESTE> \
  "<trecho do nome do evento>"
\end{verbatim}
\end{small}
\criterio{ingresso emitido; check-in ainda não validado (normal nesse ponto).}

\vspace{1mm}
\noindent\textit{Nesta tarefa o QR do participante não é lido por ninguém --- a validação é tarefa
separada do participante da Diretoria, com ingressos à parte já pré-emitidos pelo moderador.}

\section*{Fluxo 3 --- Pagamento de mensalidade (RF05, RF06, RF07, RF19)}

\estado{conta de teste do kit, com mensalidade em aberto.}

\noindent\colorbox{lightbg}{%
  \begin{minipage}{0.95\linewidth}
  \vspace{1mm}
  \textbf{Antes de começar:} avisar o participante de que o pagamento online usa uma porta de
  pagamento \textbf{simulada} (RF19) --- a tarefa valida a interação com a interface e o registro
  no sistema, não um pagamento financeiro real.
  \vspace{1mm}
  \end{minipage}%
}

\begin{enumerate}[leftmargin=6mm]
  \item Abrir a aba \textbf{Mensalidade}, localizar a cobrança em aberto do mês.
  \item Tocar \textbf{Pagar} e seguir o fluxo de pagamento online (simulado).
  \item Confirmar que o histórico passa a mostrar a competência como paga.
\end{enumerate}

\textbf{Verificação objetiva:}
\begin{verbatim}
node scripts/verificar-tarefas.js 3 <CPF-DA-CONTA-DE-TESTE>
\end{verbatim}
\criterio{mensalidade paga, com data/forma de pagamento preenchidas.}

\vfill
\noindent\footnotesize Se alguma tarefa travar por mais de \textasciitilde2 min, aplicar 1
intervenção pontual e registrar; se mesmo assim não avançar, parar e marcar como \textit{não
concluída}, seguindo para a próxima tarefa sem refazer a anterior --- protocolo completo no
\texttt{roteiro-moderador.pdf}. Depois do Fluxo 3, aplicar o questionário SUS (app) e a entrevista
semiestruturada.

\end{document}
